% STATUS: written 2026-11-05 (Guo and Li 2014, Li and Wei 2018, Ortega 2006 read in full)
% Sources read in full: Guo and Li 2014 (pp. 545--572); Li and Wei 2018 (pp. 79--92, sections 2 to 4);
% Ortega 2006 (section 1, pp. 225--232).
% Origins: path algebra and the characterization of its derivations, Guo and Li 2014
% (definition p. 547, Theorem 2.3 p. 548, Theorem 3.2 p. 554, canonical basis Theorem 3.4 p. 555);
% quiver with relations and one-point extension, Li and Wei 2018 (pp. 81--83, Theorem 3.4 p. 85);
% presentation of an incidence algebra as a path algebra with relations, Ortega 2006 (p. 232).
\section{Path algebras of quivers}\label{sec:path-incidence:path}

The triangular algebra of Chapter~\ref{ch:matrix-triangular} is a path algebra.
Take a quiver that is a line of arrows, and the algebra of its paths is the algebra
of upper triangular matrices \cite[Example~2.1, p.~82]{LiWei2018}. A quiver keeps
the order of the matrix algebra and allows branching \cite[p.~81]{LiWei2018}. This
section gives the
derivations of a path algebra, first for a plain quiver and then for a quiver with
relations.

A quiver is a directed graph with a set of vertices and a set of arrows between them
\cite[p.~547]{GuoLi2014}. Its paths form a basis of an algebra. The product of two
paths is their composition when the head of the first is the tail of the second, and
zero otherwise \cite[p.~547]{GuoLi2014}.

\begin{definition}[{\cite[p.~547]{GuoLi2014}}]\label{def:path:algebra}
Let $\Gamma=(V,E,t,h)$ be a quiver with vertex set $V$, arrow set $E$, and maps
$t,h\colon E\to V$ giving the tail and the head of an arrow. Let $P$ be the set of
paths, a path being a sequence of arrows in which the head of each arrow is the tail
of the next. The \emph{path algebra} is
\begin{equation*}
k\Gamma=\bigoplus_{p\in P}kp,\qquad
p\cdot q=\delta_{h(p),t(q)}\,pq,
\end{equation*}
where $\delta_{h(p),t(q)}$ is the delta function and $pq=0$ when $h(p)\neq t(q)$.
\end{definition}

Guo and Li characterize the derivations of $k\Gamma$ by conditions on the
coefficients \cite[p.~548]{GuoLi2014}. Write $D(p)=\sum_{q\in P}c^{p}_{q}q$ for a
linear operator $D$, with $c^{p}_{q}\in k$ \cite[p.~548]{GuoLi2014}. Call two paths
$p$ and $q$ parallel, written $p\parallel q$, when they have the same tail and the
same head \cite[p.~548]{GuoLi2014}. A path is acyclic when its head differs from its
tail, and $P_{A}$ denotes the set of acyclic paths \cite[p.~548]{GuoLi2014}.

\begin{theorem}[{\cite[Theorem~2.3, pp.~548--549]{GuoLi2014}}]\label{thm:path:derivation}
Let $\Gamma$ be a quiver and $D\colon k\Gamma\to k\Gamma$ a linear operator. Then
$D$ is a derivation if and only if it satisfies the following conditions.
\begin{enumerate}[label=(\alph*)]
\item For $v\in V$,
\begin{equation*}
D(v)=\sum_{q\in P_{A},\,t(q)=v}c^{v}_{q}q+\sum_{q\in P_{A},\,h(q)=v}c^{v}_{q}q
=\sum_{q\in P_{A},\,t(q)=v\ \text{or}\ h(q)=v}c^{v}_{q}q .
\end{equation*}
\item For $p\in P\setminus V$,
\begin{equation*}
D(p)=\sum_{q\in P_{A},\,h(q)=t(p)}c^{t(p)}_{q}qp
+\sum_{q\parallel p}c^{p}_{q}q
+\sum_{q\in P_{A},\,t(q)=h(p)}c^{h(p)}_{q}pq,
\end{equation*}
where the coefficients obey two rules.
\begin{enumerate}[label=(\roman*)]
\item For every $q\in P_{A}$, $c^{h(q)}_{q}+c^{t(q)}_{q}=0$.
\item For every $p=p_{1}p_{2}$ with $p_{1},p_{2}\in P\setminus V$ and every
$q\parallel p$,
\begin{equation*}
c^{p}_{q}=c^{p_{1}p_{2}}_{q}=
\begin{cases}
c^{p_{1}}_{qp_{1}}+c^{p_{2}}_{q_{2}},
  & p_{2}\le_{h}q\text{ with }q=q_{1}p_{2}\text{ and }p_{1}\le_{t}q\text{ with }q=p_{1}q_{2},\\
c^{p_{1}}_{q_{1}},
  & p_{2}\le_{h}q\text{ with }q=q_{1}p_{2}\text{ and }q_{1}\not\le_{t}p_{1},\\
c^{p_{2}}_{q_{2}},
  & p_{1}\le_{t}q\text{ with }q=p_{1}q_{2}\text{ and }q_{2}\not\le_{h}p_{2},\\
0, & p_{2}\not\le_{h}q\text{ and }p_{1}\not\le_{t}q,
\end{cases}
\end{equation*}
where $p\le_{t}q$ means that $q=pr$ and $p\le_{h}q$ means that $q=rp$ for some
path $r$.
\end{enumerate}
\end{enumerate}
\end{theorem}

The conditions of Theorem~\ref{thm:path:derivation} come in two families. One family
fixes the value on a vertex by values on acyclic paths at that vertex
\cite[p.~549]{GuoLi2014}. The other fixes the value on a longer path by the values
on its head, its tail and its parallel paths \cite[p.~549]{GuoLi2014}. A derivation
is therefore built from one arrow to a parallel path
\cite[pp.~553--554]{GuoLi2014}.

Fix an arrow $r\in E$ and a path $s$ parallel to $r$. Guo and Li define an operator
$D_{r,s}$ by $D_{r,s}(p)=0$ for $\ell(p)=0$, where $\ell$ is the length, and on a path
$p=p_{1}\tilde p$ by
\begin{equation*}
D_{r,s}(p)=
\begin{cases}
s\tilde p+p_{1}D_{r,s}(\tilde p), & p_{1}=r,\\
p_{1}D_{r,s}(\tilde p), & p_{1}\neq r,
\end{cases}
\end{equation*}
\cite[pp.~553--554]{GuoLi2014}. This operator is a derivation
\cite[p.~554]{GuoLi2014}.

\begin{theorem}[{\cite[Theorem~3.2, p.~554]{GuoLi2014}}]\label{thm:path:drs}
For $r\in E$ and $s\parallel r$, the linear operator
$D_{r,s}\colon k\Gamma\to k\Gamma$ defined as above is a derivation.
\end{theorem}

A quiver without arrows has no derivation but zero \cite[p.~554]{GuoLi2014}, so the
condition on the quiver cannot be dropped.

\begin{corollary}[{\cite[Corollary~3.3, p.~554]{GuoLi2014}}]\label{cor:path:existence}
There is a nonzero derivation on the path algebra $k\Gamma$ of a quiver $\Gamma$ if
and only if $\Gamma$ is a non-trivial quiver, that is, $\Gamma$ has at least one
arrow. Equivalently, $\mathrm{Der}(k\Gamma)$ is a non-zero Lie algebra if and only
if $\Gamma$ is a non-trivial quiver.
\end{corollary}

Inside the Lie algebra of derivations, the operators of Theorem~\ref{thm:path:drs}
and the inner derivations form a basis. For a path $s$, write $D_{s}(q)=sq-qs$ for
the inner derivation with respect to $s$ \cite[p.~555]{GuoLi2014}.

\begin{theorem}[{\cite[Theorem~3.4, p.~555]{GuoLi2014}}]\label{thm:path:basis}
Let $\Gamma$ be a quiver. A basis of the Lie algebra $\mathrm{Der}(k\Gamma)$ is
given by the set
\begin{equation*}
\mathcal B:=\mathcal B_{1}\cup\mathcal B_{2},\qquad
\mathcal B_{1}:=\{D_{s}\mid s\in P_{A}\},\quad
\mathcal B_{2}:=\{D_{r,s}\mid r\in E,\ s\parallel r\}.
\end{equation*}
Thus $\mathrm{Der}(k\Gamma)=D_{1}\oplus D_{2}$, where $D_{i}$ is the $k$-linear
space with basis $\mathcal B_{i}$ for $i=1,2$.
\end{theorem}

A relation on a quiver is a linear combination of parallel paths
\cite[p.~81]{LiWei2018}. With a set $\rho$ of relations, the algebra
$K(\Gamma,\rho)=K\Gamma/\langle\rho\rangle$ is a quiver with relations
\cite[p.~81]{LiWei2018}. Every basic finite-dimensional algebra over an
algebraically closed field has this form \cite[p.~81]{LiWei2018}.

For a quiver without oriented cycles, Li and Wei remove a source, and the algebra
becomes a triangular algebra over the part that remains
\cite[p.~83]{LiWei2018}. Repeating the step reaches the field, and the step is
called a one-point extension \cite[p.~83]{LiWei2018}. The step carries a statement
about the field to every finite-dimensional path algebra without oriented cycles
\cite[p.~85]{LiWei2018}.

\begin{theorem}[{\cite[Theorem~3.4, p.~85]{LiWei2018}}]\label{thm:path:jordan}
Let $K$ be a field with $\mathrm{Char}\,K\neq 2$ and $\Lambda=K(\Gamma,\rho)$ a
finite-dimensional path algebra over $K$ of the quiver $(\Gamma,\rho)$ with
relations. Then every Jordan derivation on $\Lambda$ is a derivation.
\end{theorem}

The theorem carries over to one step of that construction
\cite[p.~86]{LiWei2018}.

\begin{corollary}[{\cite[Corollary~3.6, p.~86]{LiWei2018}}]\label{cor:path:one-point}
Let $K$ be a field and $A$ a finite-dimensional $K$-algebra with every Jordan
derivation being a derivation. Let $\Lambda$ be a one-point extension of $A$ by the
bimodule $_{A}M_{K}$. Then every Jordan derivation of $\Lambda$ is a derivation.
\end{corollary}

A local derivation of a path algebra turns out to need no such freedom. Adrabi,
Bennis and Fahid prove that every local derivation of a path algebra is a
derivation, and the same for local generalized derivations
\cite[Theorem~3, p.~391]{AdrabiBennisFahid2023Local}. Their theorem settles the
Kadison question of Section~\ref{sec:local-characterizations:local} for this
whole family at once. The Lie generalized derivations of a bound quiver algebra
are treated in \cite[p.~1950]{AdrabiBennisFahid2021Lie}, and the generalized
$\sigma$-derivations of that setting have a categorical description in
\cite[p.~4754]{AdrabiBennisFahid2023Categorical}.

Ortega presents the incidence algebra of a poset as a path algebra with relations
\cite[p.~232]{Ortega2006}. Take the poset as the vertex set and put one arrow from $v$ to $w$ when $v<w$ and no
element lies between them \cite[p.~232]{Ortega2006}. The path algebra modulo the
differences of parallel paths is the incidence algebra of the poset
\cite[p.~232]{Ortega2006}. The ideal of
differences vanishes exactly when the poset is a tree, and then the incidence algebra
is hereditary \cite[p.~232]{Ortega2006}. Incidence algebras deserve their own
treatment, and Section~\ref{sec:path-incidence:incidence} gives it.

For a field $k$ and a finite simply laced quiver $Q$ without oriented cycles, Hou proves that every Lie derivation of the generalized one-point extension of the path algebra $kQ$ has the standard form with a unique standard decomposition, computes the $k$-dimensions of all its Hochschild cohomology groups for an $A$-type quiver, and describes the cup product in the cohomology ring~\cite{WlLieDerivationAndHochschildCoho2016}.
